\subsection{Affine Flats, Hulls, and Incidence}

% **Role:** incidence geometry of linear spaces without a distinguished origin.
%
% Affine flats, affine combinations, affine hulls, and flat lattices.

Throughout, $\mathcal{V}$ is finite-dimensional over a fixed field $\mathbb{F}$; affine geometry treats translated subspaces without choosing an origin.

\begin{definition} \label{def:flat-affine-geometry} \hlt{Flats}\\
Let $\mathcal{S}\leq\mathcal{V}$ and $\mathbf{x}\in\mathcal{V}$, and put $\mathbf{x}+\mathcal{S}\defeq\{\mathbf{x}+\mathbf{s}:\mathbf{s}\in\mathcal{S}\}$.
\begin{enumerate}[label=(\roman*)]
\setlength{\itemsep}{0pt}
\item The affine subset $\mathbf{x}+\mathcal{S}$ from Definition \ref{def:affine-subset-parallel-subspace} is a \hlt{flat} in $\mathcal{V}$; $\mathcal{S}$ is its \hlt{base}, and $\mathbf{x}$ is a \hlt{flat representative}.
\item The \hlt{affine geometry} on $\mathcal{V}$ is $\mathcal{A}(\mathcal{V})\defeq\{\mathbf{y}+\mathcal{T}:\mathbf{y}\in\mathcal{V},\ \mathcal{T}\leq\mathcal{V}\}$, with $\dim(\mathcal{A}(\mathcal{V}))\defeq\dim(\mathcal{V})$.
\item The \hlt{dimension} of $\mathbf{x}+\mathcal{S}$ is $\dim(\mathcal{S})$.\\
A $0$-flat is a \hlt{point}, a $1$-flat is a \hlt{line}, a $2$-flat is a \hlt{plane}, and a codimension-one flat is a \hlt{hyperplane}.
\end{enumerate}
\end{definition}

\begin{proposition} \label{prop:flat-base-and-representatives} \hlt{Flat Bases and Representatives}\\
Let $\mathcal{S},\mathcal{T}\leq\mathcal{V}$ and $\mathbf{x},\mathbf{y}\in\mathcal{V}$.
\begin{enumerate}[label=(\roman*)]
\setlength{\itemsep}{0pt}
\item $\mathbf{x}+\mathcal{S}=\mathbf{y}+\mathcal{T}\Rightarrow\mathcal{S}=\mathcal{T}$.
\item If $X=\mathbf{x}+\mathcal{S}$ and $\mathbf{z}\in X$, then $X=\mathbf{z}+\mathcal{S}$.
\end{enumerate}
\begin{proof}
\begin{enumerate}[label=(\roman*)]
\setlength{\itemsep}{0pt}
\item Suppose $\mathbf{x}+\mathcal{S}=\mathbf{y}+\mathcal{T}$ and put $\mathbf{t}_0=\mathbf{x}-\mathbf{y}\in\mathcal{T}$. Then $(\forall\,\mathbf{s}\in\mathcal{S})\,\mathbf{t}_0+\mathbf{s}=(\mathbf{x}+\mathbf{s})-\mathbf{y}\in\mathcal{T}$, so $\mathcal{S}\subseteq\mathcal{T}$.\\
Interchanging the representations proves $\mathcal{T}\subseteq\mathcal{S}$.
\item Write $\mathbf{z}=\mathbf{x}+\mathbf{s}_0$ with $\mathbf{s}_0\in\mathcal{S}$. Then $\mathbf{z}+\mathcal{S}=\mathbf{x}+(\mathbf{s}_0+\mathcal{S})=\mathbf{x}+\mathcal{S}=X$.
\end{enumerate}
\hlt{Structural Commentary}: Identify directions by two inclusions, then absorb a base displacement into the representative.\\
\hlt{Key Insight}: A flat has one base but many representatives.
\end{proof}
\end{proposition}

\begin{definition} \label{def:parallel-flats} \hlt{Parallel Flats}\\
For flats $X=\mathbf{x}+\mathcal{S}$ and $Y=\mathbf{y}+\mathcal{T}$, define $X\parallel Y\iff(\mathcal{S}\subseteq\mathcal{T}\lor\mathcal{T}\subseteq\mathcal{S})$.\\
The relation is representation-independent by Proposition \ref{prop:flat-base-and-representatives}(i).
\end{definition}

\begin{proposition} \label{prop:basic-flat-incidence} \hlt{Basic Incidence of Flats}\\
Let $X=\mathbf{x}+\mathcal{S}$ and $Y=\mathbf{y}+\mathcal{T}$ be flats in $\mathcal{V}$, and write $\mathbf{w}+X\defeq\{\mathbf{w}+\mathbf{u}:\mathbf{u}\in X\}$.
\begin{enumerate}[label=(\roman*)]
\setlength{\itemsep}{0pt}
\item $(\exists\,\mathbf{w}\in\mathcal{V})(\mathbf{w}+X\subseteq Y)\iff\mathcal{S}\subseteq\mathcal{T}$.
\item $(\exists\,\mathbf{w}\in\mathcal{V})(\mathbf{w}+X=Y)\iff\mathcal{S}=\mathcal{T}$.
\item $X\cap Y\neq\emptyset\Rightarrow(X\subseteq Y\iff\mathcal{S}\subseteq\mathcal{T})$.
\item $X\cap Y\neq\emptyset\Rightarrow(X=Y\iff\mathcal{S}=\mathcal{T})$.
\item $X\parallel Y\Rightarrow(X\subseteq Y\lor Y\subseteq X\lor X\cap Y=\emptyset)$.
\item $X\parallel Y\iff[(\exists\,\mathbf{w}\in\mathcal{V})(\mathbf{w}+X\subseteq Y)]\lor[(\exists\,\mathbf{w}\in\mathcal{V})(\mathbf{w}+Y\subseteq X)]$.
\end{enumerate}
\begin{proof}
\begin{enumerate}[label=(\roman*)]
\setlength{\itemsep}{0pt}
\item[$(i)$] ($\Rightarrow$) Fix $\mathbf{w}$ with $\mathbf{w}+X\subseteq Y$ and choose $\mathbf{z}\in\mathbf{w}+X$. Then $\mathbf{w}+X=\mathbf{z}+\mathcal{S}$ and $Y=\mathbf{z}+\mathcal{T}$ by Proposition \ref{prop:flat-base-and-representatives}(ii).\\
Hence $(\forall\,\mathbf{s}\in\mathcal{S})\,\mathbf{z}+\mathbf{s}\in\mathbf{z}+\mathcal{T}$, so $\mathcal{S}\subseteq\mathcal{T}$.
\item[$(i)$] ($\Leftarrow$) If $\mathcal{S}\subseteq\mathcal{T}$, then $(\mathbf{y}-\mathbf{x})+X=\mathbf{y}+\mathcal{S}\subseteq Y$.
\item[$(ii)$] ($\Rightarrow$) If $\mathbf{w}+X=Y$, then $\mathcal{S}=\mathcal{T}$ by Proposition \ref{prop:flat-base-and-representatives}(i).
\item[$(ii)$] ($\Leftarrow$) If $\mathcal{S}=\mathcal{T}$, then $(\mathbf{y}-\mathbf{x})+X=Y$.
\item[$(iii)$] ($\Rightarrow$) Fix $\mathbf{z}\in X\cap Y$. Then $X=\mathbf{z}+\mathcal{S}$ and $Y=\mathbf{z}+\mathcal{T}$ by Proposition \ref{prop:flat-base-and-representatives}(ii). Thus $X\subseteq Y\Rightarrow(\forall\,\mathbf{s}\in\mathcal{S})\,\mathbf{z}+\mathbf{s}\in\mathbf{z}+\mathcal{T}$, so $\mathcal{S}\subseteq\mathcal{T}$.
\item[$(iii)$] ($\Leftarrow$) Fix $\mathbf{z}\in X\cap Y$. By Proposition \ref{prop:flat-base-and-representatives}(ii), $X=\mathbf{z}+\mathcal{S}\subseteq\mathbf{z}+\mathcal{T}=Y$ whenever $\mathcal{S}\subseteq\mathcal{T}$.
\item[$(iv)$] ($\Rightarrow$) If $X=Y$, then $\mathcal{S}=\mathcal{T}$ by Proposition \ref{prop:flat-base-and-representatives}(i).
\item[$(iv)$] ($\Leftarrow$) Fix $\mathbf{z}\in X\cap Y$. If $\mathcal{S}=\mathcal{T}$, then $X=\mathbf{z}+\mathcal{S}=\mathbf{z}+\mathcal{T}=Y$ by Proposition \ref{prop:flat-base-and-representatives}(ii).
\item[$(v)$] By Definition \ref{def:parallel-flats}, $\mathcal{S}\subseteq\mathcal{T}$ or $\mathcal{T}\subseteq\mathcal{S}$. If $X\cap Y\neq\emptyset$, then $X\subseteq Y$ or $Y\subseteq X$ by (iii); otherwise $X\cap Y=\emptyset$.
\item[$(vi)$] ($\Rightarrow$) By Definition \ref{def:parallel-flats} and (i), a translate of $X$ lies in $Y$ or a translate of $Y$ lies in $X$.
\item[$(vi)$] ($\Leftarrow$) By (i), either translate-containment condition implies $\mathcal{S}\subseteq\mathcal{T}$ or $\mathcal{T}\subseteq\mathcal{S}$, so $X\parallel Y$ by Definition \ref{def:parallel-flats}.
\end{enumerate}
\hlt{Structural Commentary}: Translate incidence to base comparison in (i)--(iv), then recover parallelism in (v)--(vi).\\
\hlt{Key Insight}: A common representative converts affine incidence into subspace incidence.
\end{proof}
\end{proposition}

\begin{definition} \label{def:affine-combination-and-closed-set} \hlt{Affine Combinations}
\begin{enumerate}[label=(\roman*)]
\setlength{\itemsep}{0pt}
\item Let $n\in\mathbb{Z}_{>0}$, $\mathbf{x}_1,\ldots,\mathbf{x}_n\in\mathcal{V}$, and $r_1,\ldots,r_n\in\mathbb{F}$. The vector $\sum\limits_{j=1}^n r_j\mathbf{x}_j$ is an \hlt{affine combination} iff $\sum\limits_{j=1}^n r_j=1$.
\item A nonempty $E\subseteq\mathcal{V}$ is \hlt{affine closed} iff $(\forall\,n\in\mathbb{Z}_{>0})(\forall\,\mathbf{x}_1,\ldots,\mathbf{x}_n\in E)(\forall\,r_1,\ldots,r_n\in\mathbb{F})\,[\sum\limits_{j=1}^n r_j=1\Rightarrow\sum\limits_{j=1}^n r_j\mathbf{x}_j\in E]$.
\item A nonempty $E\subseteq\mathcal{V}$ is \hlt{two-affine closed} iff $(\forall\,\mathbf{x},\mathbf{y}\in E)(\forall\,r\in\mathbb{F})\,r\mathbf{x}+(1-r)\mathbf{y}\in E$.
\item If $\mathbf{x}\neq\mathbf{y}$, their affine line is $\overline{\mathbf{x}\mathbf{y}}\defeq\{r\mathbf{x}+(1-r)\mathbf{y}:r\in\mathbb{F}\}=\mathbf{y}+\operatorname{span}(\mathbf{x}-\mathbf{y})$.\\
Thus a nonempty $E\subseteq\mathcal{V}$ is two-affine closed iff $(\forall\,\mathbf{x},\mathbf{y}\in E)\,[\mathbf{x}\neq\mathbf{y}\Rightarrow\overline{\mathbf{x}\mathbf{y}}\subseteq E]$.
\end{enumerate}
\end{definition}

\begin{theorem} \label{thm:affine-closed-two-affine-closed} \hlt{Two-Point Affine Closure}\\
Let $\emptyset\neq E\subseteq\mathcal{V}$ and assume $\operatorname{char}(\mathbb{F})\neq2$. Then $E$ is affine closed iff $E$ is two-affine closed.
\begin{proof}
\begin{enumerate}[label=\roman*.]
\setlength{\itemsep}{0pt}
\item[$(\Rightarrow)$] Two-affine closure is the $n=2$ instance of affine closure.
\item[$(\Leftarrow)$] Assume that $E$ is two-affine closed, fix $\mathbf{x}_0\in E$, and put $\mathcal{S}=E-\mathbf{x}_0$.\\
Fix $\mathbf{s},\mathbf{t}\in\mathcal{S}$ and $c\in\mathbb{F}$. Put $\mathbf{a}\defeq\mathbf{x}_0+c\mathbf{s}=c(\mathbf{x}_0+\mathbf{s})+(1-c)\mathbf{x}_0$; then $\mathbf{a}\in E$ by two-affine closure.\\
Put $\mathbf{b}\defeq\frac12\mathbf{a}+\frac12(\mathbf{x}_0+\mathbf{t})=\mathbf{x}_0+\frac12(c\mathbf{s}+\mathbf{t})$; then $\mathbf{b}\in E$ by two-affine closure.\\
The identity $2\mathbf{b}-\mathbf{x}_0=\mathbf{x}_0+c\mathbf{s}+\mathbf{t}\in E$ holds by two-affine closure with coefficient $2$.\\
Hence $(\forall\,\mathbf{s},\mathbf{t}\in\mathcal{S})(\forall\,c\in\mathbb{F})\,c\mathbf{s}+\mathbf{t}\in\mathcal{S}$, so $\mathcal{S}\leq\mathcal{V}$ by Theorem \ref{thm:subspacetest}. Fix $n\in\mathbb{Z}_{>0}$, $\mathbf{x}_j=\mathbf{x}_0+\mathbf{s}_j\in E$, and $r_j\in\mathbb{F}$ with $\sum\limits_{j=1}^n r_j=1$.\\
Then $\sum\limits_{j=1}^n r_j\mathbf{x}_j=\mathbf{x}_0+\sum\limits_{j=1}^n r_j\mathbf{s}_j\in\mathbf{x}_0+\mathcal{S}=E$. Thus $E$ is affine closed.
\end{enumerate}
\hlt{Structural Commentary}: Translate one point to zero, derive subspace closure, then reconstruct every affine combination.\\
\hlt{Key Insight}: Midpoints and dilation by $2$ recover vector addition.
\end{proof}
\end{theorem}

\begin{remark} \label{rem:characteristic-two-affine-closure-warning} \hlt{Characteristic Two Warning}\\
The characteristic restriction is necessary. In $\mathbb{F}_2^2$, put $E\defeq\{(0,0),(1,0),(0,1)\}$.\\
Since $r\in\{0,1\}$, every two-point affine combination equals one of its endpoints, so $E$ is two-affine closed.\\
However, $(1,1)=(0,0)+(1,0)+(0,1)\notin E$ is an affine combination because $1+1+1=1$ in $\mathbb{F}_2$.
\end{remark}

\begin{theorem} \label{thm:flats-are-affine-closed} \hlt{Flats Are the Affine-Closed Sets}\\
Let $\emptyset\neq E\subseteq\mathcal{V}$. Then $E\in\mathcal{A}(\mathcal{V})$ iff $E$ is affine closed.
\begin{proof}
\begin{enumerate}[label=\roman*.]
\setlength{\itemsep}{0pt}
\item[$(\Rightarrow)$] Write $E=\mathbf{x}+\mathcal{S}$ with $\mathcal{S}\leq\mathcal{V}$. Fix $n\in\mathbb{Z}_{>0}$, $\mathbf{x}_j=\mathbf{x}+\mathbf{s}_j\in E$, and $r_j\in\mathbb{F}$ with $\sum\limits_{j=1}^n r_j=1$.\\
Then $\sum\limits_{j=1}^n r_j\mathbf{x}_j=\mathbf{x}+\sum\limits_{j=1}^n r_j\mathbf{s}_j\in E$, so $E$ is affine closed.
\item[$(\Leftarrow)$] Assume that $E$ is affine closed, fix $\mathbf{x}\in E$, and put $\mathcal{S}=E-\mathbf{x}$. Fix $\mathbf{s}_1,\mathbf{s}_2\in\mathcal{S}$ and $c\in\mathbb{F}$, and choose $\mathbf{x}_1,\mathbf{x}_2\in E$ with $\mathbf{s}_j=\mathbf{x}_j-\mathbf{x}$.\\
Then $c\mathbf{s}_1+\mathbf{s}_2+\mathbf{x}=c\mathbf{x}_1+\mathbf{x}_2-c\mathbf{x}\in E$ by affine closure, so $c\mathbf{s}_1+\mathbf{s}_2\in\mathcal{S}$. Hence $\mathcal{S}\leq\mathcal{V}$ by Theorem \ref{thm:subspacetest}, and $E=\mathbf{x}+\mathcal{S}\in\mathcal{A}(\mathcal{V})$.
\end{enumerate}
\hlt{Structural Commentary}: Check translated-subspace closure forward, then recover a base by shifting one point to zero.\\
\hlt{Key Insight}: Affine closure becomes linear closure after translation.
\end{proof}
\end{theorem}

\begin{theorem} \label{thm:flats-are-two-affine-closed} \hlt{Flats Are the Two-Affine-Closed Sets}\\
Let $\emptyset\neq E\subseteq\mathcal{V}$ and assume $\operatorname{char}(\mathbb{F})\neq2$. Then $E\in\mathcal{A}(\mathcal{V})$ iff $E$ is two-affine closed.
\begin{proof}
\begin{enumerate}[label=\roman*.]
\setlength{\itemsep}{0pt}
\item[$(\Rightarrow)$] If $E\in\mathcal{A}(\mathcal{V})$, then $E$ is affine closed by Theorem \ref{thm:flats-are-affine-closed} and two-affine closed by Theorem \ref{thm:affine-closed-two-affine-closed}.
\item[$(\Leftarrow)$] If $E$ is two-affine closed, then $E$ is affine closed by Theorem \ref{thm:affine-closed-two-affine-closed} and $E\in\mathcal{A}(\mathcal{V})$ by Theorem \ref{thm:flats-are-affine-closed}.
\end{enumerate}
\hlt{Structural Commentary}: Compose the two preceding characterizations in each direction.\\
\hlt{Key Insight}: Binary affine closure detects translated subspaces when $2$ is invertible.
\end{proof}
\end{theorem}

\begin{definition} \label{def:affine-hull-and-affine-span} \hlt{Affine Hull and Affine Span}\\
Let $\emptyset\neq E\subseteq\mathcal{V}$.
\begin{enumerate}[label=(\roman*)]
\setlength{\itemsep}{0pt}
\item The \hlt{affine hull} is $\operatorname{affhull}(E)\defeq\bigcap\limits\{X\in\mathcal{A}(\mathcal{V}):E\subseteq X\}$; the defining family contains $\mathcal{V}$.
\item The \hlt{affine span} is $\operatorname{affspan}(E)\defeq\{\sum\limits_{j=1}^n r_j\mathbf{e}_j:n\in\mathbb{Z}_{>0},\ \mathbf{e}_j\in E,\ r_j\in\mathbb{F},\ \sum\limits_{j=1}^n r_j=1\}$.
\end{enumerate}
\end{definition}

\begin{theorem} \label{thm:affine-hull-span-formula} \hlt{Affine Hull Formula}\\
Let $\emptyset\neq E\subseteq\mathcal{V}$, fix $\mathbf{x}\in E$, and put $\mathcal{S}_{\mathbf{x}}\defeq\operatorname{span}(E-\mathbf{x})$.
\begin{enumerate}[label=(\roman*)]
\setlength{\itemsep}{0pt}
\item $\operatorname{affhull}(E)=\operatorname{affspan}(E)=\mathbf{x}+\mathcal{S}_{\mathbf{x}}$.
\item $\dim(\operatorname{affhull}(E))=\dim(\mathcal{S}_{\mathbf{x}})$.
\end{enumerate}
\begin{proof}
\begin{enumerate}[label=(\roman*)]
\setlength{\itemsep}{0pt}
\item If $\mathbf{z}=\sum\limits_{j=1}^n r_j\mathbf{e}_j\in\operatorname{affspan}(E)$, then $\mathbf{z}=\mathbf{x}+\sum\limits_{j=1}^n r_j(\mathbf{e}_j-\mathbf{x})\in\mathbf{x}+\mathcal{S}_{\mathbf{x}}$.\\
Conversely, if $\mathbf{z}=\mathbf{x}+\sum\limits_{j=1}^n a_j(\mathbf{e}_j-\mathbf{x})\in\mathbf{x}+\mathcal{S}_{\mathbf{x}}$, then $\mathbf{z}=(1-\sum\limits_{j=1}^n a_j)\mathbf{x}+\sum\limits_{j=1}^n a_j\mathbf{e}_j\in\operatorname{affspan}(E)$.\\
Hence $\operatorname{affspan}(E)=\mathbf{x}+\mathcal{S}_{\mathbf{x}}$, a flat containing $E$.\\
Every flat containing $E$ contains $\operatorname{affspan}(E)$ by Theorem \ref{thm:flats-are-affine-closed}, so this flat equals $\operatorname{affhull}(E)$ by Definition \ref{def:affine-hull-and-affine-span}(i).
\item By (i) and Definition \ref{def:flat-affine-geometry}(iii), $\dim(\operatorname{affhull}(E))=\dim(\mathbf{x}+\mathcal{S}_{\mathbf{x}})=\dim(\mathcal{S}_{\mathbf{x}})$.
\end{enumerate}
\hlt{Structural Commentary}: Prove the translated-span formula by two inclusions, then read dimension from its base.\\
\hlt{Key Insight}: Differences from one point encode every affine direction.
\end{proof}
\end{theorem}

\begin{lemma} \label{lem:intersection-of-flats} \hlt{Intersections of Flats}\\
Let $\emptyset\neq I$ and let $X_i=\mathbf{x}_i+\mathcal{S}_i$ be flats indexed by $i\in I$. If $\mathbf{z}\in\bigcap\limits_{i\in I}X_i$, then $\bigcap\limits_{i\in I}X_i=\mathbf{z}+\bigcap\limits_{i\in I}\mathcal{S}_i$.
\begin{proof}
By Proposition \ref{prop:flat-base-and-representatives}(ii), $(\forall\,i\in I)\,X_i=\mathbf{z}+\mathcal{S}_i$. Hence $(\forall\,\mathbf{u}\in\mathcal{V})\,[\mathbf{u}\in\bigcap\limits_{i\in I}X_i\iff\mathbf{u}-\mathbf{z}\in\bigcap\limits_{i\in I}\mathcal{S}_i]$, which proves the formula.\\
\hlt{Structural Commentary}: Replace every representative by the common point, then intersect the remaining bases.\\
\hlt{Key Insight}: Common incidence synchronizes all affine origins.
\end{proof}
\end{lemma}

\begin{theorem} \label{thm:lattice-of-flats} \hlt{Lattice of Flats}\\
Let $\mathcal{A}_0(\mathcal{V})\defeq\mathcal{A}(\mathcal{V})\cup\{\emptyset\}$, ordered by inclusion.
\begin{enumerate}[label=(\roman*)]
\setlength{\itemsep}{0pt}
\item $(\forall\,\mathcal{C}\subseteq\mathcal{A}_0(\mathcal{V}))\,\bigwedge\mathcal{C}=\bigcap\limits_{X\in\mathcal{C}}X$.
\item If $\mathcal{C}\subseteq\mathcal{A}_0(\mathcal{V})$ and $U\defeq\bigcup\limits_{X\in\mathcal{C}}X$, then $U\neq\emptyset\Rightarrow\bigvee\mathcal{C}=\operatorname{affhull}(U)$, while $U=\emptyset\Rightarrow\bigvee\mathcal{C}=\emptyset$.
\item $\mathcal{A}_0(\mathcal{V})$ is a complete lattice.
\end{enumerate}
\begin{proof}
\begin{enumerate}[label=(\roman*)]
\setlength{\itemsep}{0pt}
\item Fix $\mathcal{C}\subseteq\mathcal{A}_0(\mathcal{V})$ and put $M=\bigcap\limits_{X\in\mathcal{C}}X$. If $\mathcal{C}=\emptyset$, then $M=\mathcal{V}$. If $\mathcal{C}\neq\emptyset$ and $M=\emptyset$, then $M\in\mathcal{A}_0(\mathcal{V})$.\\
If $\mathcal{C}\neq\emptyset$ and $M\neq\emptyset$, choose $\mathbf{z}\in M$ and let $\mathcal{S}_X$ be the unique base of $X$; then $M=\mathbf{z}+\bigcap\limits_{X\in\mathcal{C}}\mathcal{S}_X$ by Lemma \ref{lem:intersection-of-flats}, so $M$ is a flat.\\
Thus $M\in\mathcal{A}_0(\mathcal{V})$. It is a lower bound of $\mathcal{C}$, and any lower bound is contained in $M$, so $M=\bigwedge\mathcal{C}$.
\item Fix $\mathcal{C}\subseteq\mathcal{A}_0(\mathcal{V})$ and put $U=\bigcup\limits_{X\in\mathcal{C}}X$.\\
If $U=\emptyset$, then $\emptyset$ is the least upper bound of $\mathcal{C}$. If $U\neq\emptyset$, then $\operatorname{affhull}(U)$ is a flat containing each $X\in\mathcal{C}$ by Theorem \ref{thm:affine-hull-span-formula}.\\
If $B$ is an upper bound of $\mathcal{C}$, then $U\subseteq B$ and $\operatorname{affhull}(U)\subseteq B$, so $\operatorname{affhull}(U)=\bigvee\mathcal{C}$.
\item The arbitrary meets and joins in (i)--(ii) prove completeness by Definition \ref{def:lattice-vocabulary}(iv).
\end{enumerate}
\hlt{Structural Commentary}: Construct arbitrary meets by intersection and joins by affine hull, then invoke lattice completeness.\\
\hlt{Key Insight}: Empty incidence must be adjoined to close arbitrary meets.
\end{proof}
\end{theorem}

\begin{remark} \label{rem:empty-affine-hull-lattice-convention} \hlt{Empty Affine Hull Convention}\\
Within $\mathcal{A}_0(\mathcal{V})$, set $\operatorname{affhull}(\emptyset)\defeq\emptyset$; then Theorem \ref{thm:lattice-of-flats}(ii) has the uniform form $\bigvee\mathcal{C}=\operatorname{affhull}(\bigcup\limits_{X\in\mathcal{C}}X)$.
\end{remark}

\begin{corollary} \label{cor:join-of-two-flats-formula} \hlt{Join of Two Flats}\\
Let $X=\mathbf{x}+\mathcal{S}$ and $Y=\mathbf{y}+\mathcal{T}$ be flats.
\begin{enumerate}[label=(\roman*)]
\setlength{\itemsep}{0pt}
\item $X\vee Y=\mathbf{x}+\operatorname{span}(\{\mathbf{x}-\mathbf{y}\}\cup\mathcal{S}\cup\mathcal{T})$.
\item $X\cap Y\neq\emptyset\Rightarrow X\vee Y=\mathbf{x}+(\mathcal{S}+\mathcal{T})$.
\end{enumerate}
\begin{proof}
\begin{enumerate}[label=(\roman*)]
\setlength{\itemsep}{0pt}
\item By Theorem \ref{thm:lattice-of-flats}(ii) and Theorem \ref{thm:affine-hull-span-formula}(i), $X\vee Y=\mathbf{x}+\operatorname{span}((X\cup Y)-\mathbf{x})$.\\
Since $(X\cup Y)-\mathbf{x}=\mathcal{S}\cup((\mathbf{y}-\mathbf{x})+\mathcal{T})$, its span is $\operatorname{span}(\{\mathbf{x}-\mathbf{y}\}\cup\mathcal{S}\cup\mathcal{T})$.
\item Fix $\mathbf{z}=\mathbf{x}+\mathbf{s}=\mathbf{y}+\mathbf{t}\in X\cap Y$. Then $\mathbf{x}-\mathbf{y}=\mathbf{t}-\mathbf{s}\in\mathcal{S}+\mathcal{T}$, so the span in (i) equals $\mathcal{S}+\mathcal{T}$.
\end{enumerate}
\hlt{Structural Commentary}: Expand the affine hull around $\mathbf{x}$, then test whether the displacement adds a direction.\\
\hlt{Key Insight}: Intersection absorbs the displacement into the two bases.
\end{proof}
\end{corollary}

\begin{corollary} \label{cor:dimension-of-join-of-flats} \hlt{Dimension of the Join of Two Flats}\\
Let $X=\mathbf{x}+\mathcal{S}$ and $Y=\mathbf{y}+\mathcal{T}$ be flats.
\begin{enumerate}[label=(\roman*)]
\setlength{\itemsep}{0pt}
\item $X\cap Y\neq\emptyset\Rightarrow\dim(X\vee Y)=\dim(X)+\dim(Y)-\dim(X\cap Y)$.
\item $X\cap Y=\emptyset\Rightarrow\dim(X\vee Y)=\dim(\mathcal{S}+\mathcal{T})+1$.
\end{enumerate}
\begin{proof}
\begin{enumerate}[label=(\roman*)]
\setlength{\itemsep}{0pt}
\item Fix $\mathbf{z}\in X\cap Y$. By Corollary \ref{cor:join-of-two-flats-formula}(ii) and Lemma \ref{lem:intersection-of-flats}, the bases of $X\vee Y$ and $X\cap Y$ are $\mathcal{S}+\mathcal{T}$ and $\mathcal{S}\cap\mathcal{T}$, respectively.\\
Hence $\dim(X\vee Y)=\dim(\mathcal{S})+\dim(\mathcal{T})-\dim(\mathcal{S}\cap\mathcal{T})$ by Theorem \ref{thm:dimension-formula-for-sum}, which is the stated identity by Definition \ref{def:flat-affine-geometry}(iii).
\item By Corollary \ref{cor:join-of-two-flats-formula}(i), the base of $X\vee Y$ is $\mathcal{U}\defeq\operatorname{span}(\{\mathbf{x}-\mathbf{y}\}\cup\mathcal{S}\cup\mathcal{T})$.\\
If $\mathbf{x}-\mathbf{y}\in\mathcal{S}+\mathcal{T}$, choose $\mathbf{s}\in\mathcal{S}$ and $\mathbf{t}\in\mathcal{T}$ with $\mathbf{x}-\mathbf{y}=\mathbf{s}+\mathbf{t}$; then $\mathbf{x}-\mathbf{s}=\mathbf{y}+\mathbf{t}\in X\cap Y$, a contradiction.\\
Thus $\mathcal{U}=(\mathcal{S}+\mathcal{T})\oplus\operatorname{span}(\mathbf{x}-\mathbf{y})$, and $\dim(X\vee Y)=\dim(\mathcal{S}+\mathcal{T})+1$ by Corollary \ref{cor:dimension-direct-sum}.
\end{enumerate}
\hlt{Structural Commentary}: Apply the subspace dimension formula when flats meet; isolate one new direction otherwise.\\
\hlt{Key Insight}: Disjointness makes the representative displacement linearly independent.
\end{proof}
\end{corollary}
